\begin{afterword}
\summaryhead

The post sums up the sequence's case for many-worlds. The laws found in microscopic experiments are simple, and the simplest generalization is that they hold unchanged at every scale. Applied to Schrödinger's cat, they put the cat and the observer in superposition. Decoherence keeps the versions apart, so each observer sees one outcome. Other worlds are therefore a consequence of the laws, not an added assumption, and Occam's razor, properly understood, does not count against them.

The post grants one open puzzle: why probabilities follow the squared amplitudes, the Born rule. It reports Hanson's ``mangled worlds'' proposal, doubts the decision-theoretic derivations of Deutsch and Wallace, and hopes for an answer without new laws, while granting that one might be needed.

Against a single world it gives two reasons. No fundamental law could pick out one high-level world. And by Bell's theorem any single-world theory needs an influence faster than light, which violates special relativity. Non-realism is dismissed as a way of avoiding the question. The post grants that some effect might appear only at large scales and that its argument might contain a hidden assumption, but it rates one world about as likely as spinning black holes violating conservation of angular momentum. It blames the continuing debate on historical accident and on physicists' weak grasp of Occam's razor, and declares that many-worlds ``wins outright.''

\responsehead

Much of the post is sound, and that should be said first. It states the main weak point of its own view, the Born rule, openly, and its suspicion of the decision-theoretic answers is shared by much of the published literature. Several of its steps are standard. Specialist sources agree that many-worlds is the most economical theory if laws are counted rather than objects, and that any single-world account must be nonlocal. Its figures on the Bell experiments are correct.

The case for ``wins outright'' depends on two premises, and each claims more than the literature grants. The first is that other worlds are a ``logical consequence'' of the laws. They are, but only if the wavefunction is all there is. Bohm's theory keeps the same wave equation everywhere, adds particle positions, has one world, and predicts the same results. The second is that a single world ``would violate Special Relativity.'' It must be nonlocal; that it must violate relativity is disputed, and relativistic versions of the GRW collapse theory were published before the post. The argument does bear on Bohm-type theories, which need a preferred slicing of spacetime. So it rules out some single-world theories, not all of them.

The simplicity comparison is also less one-sided than the post's closing claims suggest. The usual many-worlds answer to the Born puzzle adds a postulate about probability. How there can be probability at all when every outcome occurs is, in the words of an encyclopedia entry that argues for many-worlds, ``arguably the most serious difficulty.''

The closing paragraphs move from these premises to ``no rational controversy'' and an ``established fact.'' They explain the remaining disagreement by historical accident, by physicists' poor grasp of Occam's razor and by fear for their careers; earlier, the post put the question of one world down to a ``sentimental attachment'' to intuition. It gives no evidence for any of these, though others have made the historical point too. Specialists who know the formal arguments have published objections on probability and on what the wavefunction can describe, and open defenders of many-worlds had written for decades before the post. The field has not moved toward the post's verdict: many-worlds was the favourite of 15 to 18 per cent of respondents in polls from 1998 to 2025.

In short: a clear statement of the real strengths of many-worlds and of its main open problem, which then treats two disputed premises as settled and puts the remaining disagreement down to the ignorance, fear and sentiment of those who disagree.
\end{afterword}
